\section{Matriz de atitude e outras composições}

Para representar a atitude (orientação) de um corpo rígido, neste caso, o manipulador robótico, é necessário detalhar a relação entre os \emph{frames (ou sistema de coordenadas)} e um sistema de referência global.

\subsection{Relação entre frames e sistema de referência global}
Como existem múltiplos
\emph{frames} envolvidos, como o Inercial (centrado na Terra, $\{I\}$);
o \emph{frame} Local da espaçonave, $\{LVLH\}$, que por sua vez está no centro de massa do \emph{chaser};
e o \emph{corpo $\{B\}$}, fixado no centro de massa do manipulador robótico,
se faz necessário definir a relação entre esses \emph{frames} e o sistema de referência global.
A matriz de rotação é utilizada para descrever as coordenadas de um vetor de um sistema para outro \cite{craig2005}.

\begin{equation}
\label{eq:LocalparaInercial}
{}^{I}\mathbf{R}_{B} = {}^{I}\mathbf{R}_{LVLH}\;{}^{LVLH}\mathbf{R}_{B}
\end{equation}

Enquanto que a matriz de rotação pode ser descrita por:
\begin{equation}
\label{eq:LocalparaInercial_yawpitchroll}
{}^{LVLH}\mathbf{R}_{B} (\psi,\theta,\phi) = \mathbf{R}_z(\psi)\,\mathbf{R}_y(\theta)\,\mathbf{R}_x(\phi)
\end{equation}

A composição de \eqref{eq:LocalparaInercial_yawpitchroll} corresponde à sequência extrínseca XYZ adotada anteriormente, e para melhor compreensão dos termos nesta seção, os símbolos de rotação em cada eixo principal foram definidos na Tabela~\ref{tab:euler_LVLH}.

\begin{table}[ht]
    \centering
    \footnotesize
    \caption{Relação entre ângulos de Euler (sequência ZYX) e eixos do sistema LVLH.}
    \label{tab:euler_LVLH}
    \begin{tabular}{lll}
        \hline
        \textbf{Ângulo} & \textbf{Eixo de rotação} & \textbf{Descrição} \\
        \hline
        $\phi$ (roll, rolagem)
        & $x$ (radial)
        & Rotação em torno do eixo radial \\

        $\theta$ (pitch, arremesso)
        & $y$ (along-track)
        & Rotação em torno do eixo tangente à órbita \\

        $\psi$ (yaw, guinada)
        & $z$ (cross-track)
        & Rotação lateral em torno do eixo normal ao plano orbital \\

        \hline
    \end{tabular}
\end{table}

Após a definição das matrizes homogêneas, introduz-se também o operador \emph{chapéu} \cite{murray1994mathematical}. Esse operador auxilia nas transformações de coordenadas e na formulação das equações dinâmicas, sendo definido como a matriz assimétrica associada a um vetor $a = [x, y, z]^T$.

\begin{equation}
    \label{eq:matriz_assimetrica}
    \hat{a} = 
    \begin{bmatrix}
        0   & -z  &  y \\
        z   &  0  & -x \\
        -y  &  x  &  0
    \end{bmatrix}
\end{equation}

Essa matriz possui a propriedade de que $\hat{a} b = a \times b$, ou seja, a multiplicação matricial equivale ao produto vetorial tradicional.
Como o manipulador robótico deste documento possui 5 juntas, pode-se adaptar \eqref{eq:produtorio_generico} para essa totalidade de juntas, como indicado a seguir:
\begin{equation}
{}^{0}\!T_E(q)=
A_1(\theta_1)
A_2(\theta_2)
A_3(\theta_3)
A_4(\theta_4)
A_5(d_5).
\end{equation}

Para extrair a orientação, utiliza-se \eqref{eq:rotacao_MatrizHomogenea_orientacao} e para extrair a posição utiliza-se \eqref{eq:posicao_MatrizHomogenea_posicao}, conforme:

\begin{equation}
    \label{eq:rotacao_MatrizHomogenea_orientacao}
\mathbf{R}_E(q) = \big[{}^{0}\mathbf{T}_E(q)\big]_{1:3,\,1:3},\qquad
\end{equation}

Onde o ponto do Efetuador $\{E\}$ é definido como:
\begin{equation}
    \label{eq:posicao_MatrizHomogenea_posicao}
\vec{p}_E(q) = \big[{}^{0}\mathbf{T}_E(q)\big]_{1:3,\,4}.
\end{equation}
